\documentclass[../main]{subfiles}
\begin{document}

\section{Realización de cálculo de recipientes sometidos a presión}

En especialmente en el diseño y mantenimiento de estructuras y sistemas de almacenamiento de fluidos, es fundamental comprender y calcular las fuerzas y presiones a las que están sometidos los recipientes. Estas presiones pueden variar según la naturaleza del fluido, la temperatura, el diseño del recipiente y otros factores. En este trabajo práctico, exploraremos algunos conceptos fundamentales y fórmulas relacionadas con el cálculo de recipientes sometidos a presión.

\subsection{Presión y fuerza}

La presión ($P$) es una magnitud física que mide la fuerza por unidad de área que ejerce un fluido sobre una superficie. Se expresa en unidades de presión como Pascal (Pa), atmosfera (atm), bar (bar), entre otros.

La fuerza ($F$) ejercida por un fluido sobre una superficie se puede calcular mediante la fórmula:

\begin{equation}
F = P \times A
\end{equation}

Donde:
\begin{itemize}
    \item $F$ es la fuerza ejercida por el fluido (N).
    \item $P$ es la presión del fluido (Pa).
    \item $A$ es el área sobre la cual actúa la presión (m$^2$).
\end{itemize}

\subsection{Ley de Pascal}

La Ley de Pascal establece que la presión aplicada a un fluido confinado se transmite sin cambios en todas las direcciones y en todos los puntos del fluido. Matemáticamente, se expresa como:

\begin{equation}
P_1 = P_2
\end{equation}

Donde:
\begin{itemize}
    \item $P_1$ es la presión aplicada en un punto del fluido (Pa).
    \item $P_2$ es la presión en otro punto del fluido (Pa).
\end{itemize}

Esta ley es fundamental para comprender cómo se distribuye la presión dentro de un recipiente.

\subsection{Ecuación de la presión hidrostática}

La presión hidrostática es la presión que ejerce un fluido en reposo debido a la fuerza de la gravedad. Se calcula mediante la ecuación:

\begin{equation}
P = \rho \times g \times h
\end{equation}

Donde:
\begin{itemize}
    \item $P$ es la presión hidrostática (Pa).
    \item $\rho$ es la densidad del fluido (kg/m$^3$).
    \item $g$ es la aceleración debido a la gravedad (m/s$^2$).
    \item $h$ es la altura del fluido sobre el punto de referencia (m).
\end{itemize}

Esta ecuación nos permite calcular la presión en cualquier punto dentro de un fluido en reposo, como un tanque de almacenamiento.

\subsection{Preguntas para investigar}

\begin{enumerate}
    \item ¿Cuáles son los principales tipos de recipientes sometidos a presión utilizados en la industria?
    \item ¿Cómo influye la temperatura en el cálculo de la presión de un recipiente?
    \item ¿Qué normativas o códigos regulan el diseño y mantenimiento de recipientes sometidos a presión?
    \item ¿Cuáles son los materiales más comunes utilizados en la construcción de recipientes a presión?
    \item ¿Qué métodos se utilizan para inspeccionar y evaluar la integridad de un recipiente sometido a presión?
    \item ¿Cómo se pueden mitigar los riesgos asociados con la operación de recipientes a presión?
    \item ¿Qué papel juegan las juntas y sellos en la prevención de fugas en recipientes a presión?
    \item ¿Cuál es la importancia del alivio de presión en el diseño de recipientes a presión?
    \item ¿Cómo se calcula el espesor requerido para un recipiente a presión según los códigos de diseño?
    \item ¿Qué consideraciones se deben tener en cuenta al realizar pruebas hidrostáticas en un recipiente a presión?
\end{enumerate}

\subsection{Problemas}

\begin{enumerate}
    \item Un tanque cilíndrico tiene un diámetro de 2 metros y contiene agua hasta una altura de 5 metros. ¿Cuál es la presión hidrostática en el fondo del tanque? (Considere la densidad del agua como 1000 kg/m$^3$ y la aceleración debido a la gravedad como 9.81 m/s$^2$).
    \item Se aplica una presión de 3 atmósferas a un tanque esférico cuyo radio es de 1 metro. ¿Cuál es la fuerza total ejercida por el fluido sobre la superficie interna del tanque?
    \item ¿Cuál es la presión máxima que puede soportar una tubería de acero con un diámetro de 10 cm y un espesor de 5 mm, si se utiliza para transportar agua a una presión de 8 bar?
    \item Se tiene un recipiente cúbico con una arista de 3 metros. Si se llena con aceite de densidad 850 kg/m$^3$ hasta una altura de 2 metros, ¿cuál es la fuerza ejercida por el aceite en una de las caras laterales del cubo?
    \item ¿Cuál debería ser el espesor mínimo de un tanque cilíndrico de acero con un diámetro de 4 metros y que operará a una presión interna de 5 bar, según la normativa ASME (Sociedad Americana de Ingenieros Mecánicos)?
    \item Un recipiente esférico tiene un radio de 2 metros y está lleno de gas a una presión de 10 bar. Si se perfora una pequeña abertura en la superficie del recipiente, ¿cuál será la velocidad inicial del gas que sale por la abertura?
    \item Se requiere almacenar 1000 litros de amoníaco líquido a una presión de 6 atmósferas. ¿Cuál debería ser el volumen mínimo de un tanque cilíndrico para cumplir con este requisito?
    \item ¿Cómo afecta la presión del fluido al diseño de las válvulas de alivio de presión en un sistema de tuberías?
    \end{enumerate}
 \end{document}
